% id: 24-20
% book: 베이직쎈 미적분1 · 01 함수의 극한
% section: 개념 12 함수의 극한의 도형에의 활용
% passage: 오른쪽 그림과 같이 함수 $y=x^2$의 그래프 위의 원점~$\pt{O}$가 아닌 점~$\pt{P}(t{,}\mkern3mu\mathopen{}t^2)$에 대하여 점~$\pt{P}$를 지나고 직선~$\pt{OP}$와 수직인 직선이 $y$축과 만나는 점의 $y$좌표를 $f(t)$라 할 때, 다음을 구하시오.
% figure: tikz:fig-24-17
% answer: 
% answer_source: 
% variant_level: 0 (원본 전사)
% 이 파일은 build.mjs 가 items.json 에서 만든다. 고칠 때는 items.json 을 고친다.
\smash{\raisebox{1.5mm}{\dmkichul{베이직쎈 미적분1 24쪽 20번}}}
점~$\pt{P}$가 함수 $y=x^2$의 그래프를 따라 원점~$\pt{O}$에 한없이 가까워질 때, $f(t)$의 값이 한없이 가까워지는 값
